\section{用向量丛解释}

\ob{2.1}{ob:2.1}
对我们来说，一个关键的观察是：集合 $\GL_n(\cK)\backslash\GL_n(\A)/\GL_n(\sO)$ 可以改用几何语言重述。

记 $\Bun_n$ 为曲线 $X$ 上秩 $n$ 向量丛的同构类集合。我们断言：$\Bun_n$ 与 $\GL_n(\cK)\backslash\GL_n(\A)/\GL_n(\sO)$ 之间存在自然双射。

事实上，给定向量丛 $M$，它在 Zariski 拓扑下局部同构于平凡丛 $\cO_X^{\oplus n}$。取它在 $X$ 的一般点处的平凡化 $M\otimes_{\cO_X}\cK\simeq\cK^n$，并在各个局部位处取平凡化 $M\otimes_{\cO_X}\cO_x\simeq\cO_x^n$。

对每个 $x\in|X|$ 复合这些同构，得到
\[
\cK_x^n\simeq\cO_x^n\otimes_{\cO_x}\cK_x\simeq M\otimes_{\cO_X}\cK_x
\simeq M\otimes_{\cO_X}\cK\otimes_{\cK}\cK_x\simeq\cK^n\otimes_{\cK}\cK_x\simeq\cK_x^n,
\]
即 $\GL_n(\A)$ 的一个元素。改变上述平凡化，效果是在左边乘 $\GL_n(\cK)$ 的元素，在右边乘 $\prod_{x\in|X|}\GL_n(\cO_x)=\GL_n(\sO)$ 的元素。

\begin{annote}
这里看似要作无穷多个选择：一处一般点的平凡化，加上每个位处的局部平凡化。实际上改变某个 $x$ 处的局部平凡化，相当于在 $\GL_n(\cK_x)$ 中右乘；所有位合起来，正是右乘 $\GL_n(\sO)$。改变一般点的平凡化则相当于左乘 $\GL_n(\cK)$。于是双射只依赖向量丛的同构类。另须注意，此处 $\Bun_n$ 还只是一个集合（同构类的集合）；\S3.1 才把它升级为模叠 $\sBun_n$，届时「向量丛的同构类」被「向量丛的范畴」取代。\par\vspace{0.15em}{\raggedleft 【来源：原书 \S2.1、\S3.1（A）】\par}
\end{annote}

\ob{2.2}{ob:2.2}
\emph{例}\quad 当 $n=1$ 时，$\GL_1(\A)/\GL_1(\sO)$ 是 $X$ 上的除子集合。因此双商 $\GL_1(\cK)\backslash\GL_1(\A)/\GL_1(\sO)$ 是除子类（即除子模主除子）的集合。上述双射重新建立了熟知的事实：除子类群等同于线丛的同构类群。

\ob{2.3}{ob:2.3}
现在说明如何在几何语言中看待 Hecke 算子。固定点 $x\in X$ 与整数 $i$，$1\le i\le n$。

记 $H_x^i$ 为下述数据的同构类集合：$(M,M',\beta)$，其中 $M$ 与 $M'$ 是 $X$ 上秩 $n$ 向量丛，$\beta$ 是 $M$ 到 $M'$ 的单射（视二者为 $X$ 上的凝聚层），使得商 $M'/M$ 长度为 $i$，且概形论地支撑在 $x$ 处。（后一条件等价于说 $M'/M$ 非典范地同构于 $k_x^i$，其中 $k_x$ 是 $x$ 处的剩余域。）

有两个自然映射，记为 $h^{\leftarrow}$ 与 $h^{\rightarrow}$，从 $H_x^i$ 到 $\Bun_n$，分别把上述三元组送到 $M$ 与 $M'$。

\begin{othm}{引理}{2.4}{lem:2.4}
给定 $M\in\Bun_n$，$h^{\leftarrow}$ 在其上的纤维可典范地等同于 $k_x$ 上向量空间 $M_x$ 中 $i$ 维子空间的集合。换言之，$(h^{\leftarrow})^{-1}(M)$ 是 Grassmann 簇 $\Gr^i(M_x)$ 的 $k_x$-点集。
\end{othm}

对 $h^{\rightarrow}$ 同样成立；此时得到 $(h^{\rightarrow})^{-1}(M')$ 是 $\Gr^{n-i}(M'_x)$ 的 $k_x$-点空间。

\begin{proof}
设 $(M,M',\beta)$ 如上。$\beta$ 虽是凝聚层的单射，取 $x$ 处的纤维后却不再单射，而得到正合列
\[
0\to\ker(\beta|_x)\to M_x\to M'_x\to(M'/M)_x\to0.
\]
特别地，$\ker(\beta|_x)$ 是 $M_x$ 的 $i$ 维子空间。

反之，给定子空间 $V\subset M_x$，定义 $M'$ 为 $M(x)$ 的子层，其局部截面形如 $t^{-1}\cdot m$，其中 $t$ 是 $X$ 在 $x$ 处的局部参数，$m$ 是 $M$ 的局部截面且在 $x$ 处的值属于 $V$。
\end{proof}

\begin{annote}
$M(x)$ 是原书的简写，指把 $M$ 沿 $x$ 扭转一次，即 $M\otimes_{\cO_X}\cO_X(x)$。它比 $M$ 多出允许在 $x$ 处有一阶极点的截面，正是形如 $t^{-1}m$ 的那些。要求 $m$ 在 $x$ 处的值落在 $V$ 内，就限定了 $M'$ 在 $x$ 处纤维的大小；这一构造的逆给出 $\ker(\beta|_x)=V$，从而正合列中的 $\ker$ 恰好跑遍 $M_x$ 的 $i$ 维子空间。\par\vspace{0.15em}{\raggedleft 【来源：原书 \S2.4（A）；$M(x)$ 的常规含义（B）】\par}
\end{annote}

\ob{2.5}{ob:2.5}
注意，由于 $k_x$ 是有限域，$h^{\leftarrow}$ 与 $h^{\rightarrow}$ 的纤维都是有限集。

考虑从空间 $\Funct(\Bun_n)$ 到自身的算子
\[
f\mapsto(h^{\rightarrow})_!((h^{\leftarrow})^*(f)).
\]
（上标 $*$ 表示沿指定态射的函数拉回，下标 $!$ 表示沿纤维求和；可与附录中的记号对照。）
\[
\Bun_n\xleftarrow{h^{\leftarrow}}H_x^i\xrightarrow{h^{\rightarrow}}\Bun_n.
\]

\begin{othm}{命题}{2.6}{prop:2.6}
上述算子等于 $T_x^i(f)$。
\end{othm}

\ob{2.7}{ob:2.7}
\emph{插叙}\quad 我们指出与模形式情形的如下类比。设 $N$ 是自然数，$p$ 是与 $N$ 互素的素数。我们将研究模曲线 $X(N)$ 上的函数（更确切地说，某些丛的截面）（它与本文的曲线 $X$ 无关）。

按定义，$X(N)$ 参数化对 $(E,\alpha)$，其中 $E$ 是（$\mathbb{C}$ 上的）椭圆曲线，$\alpha$ 是 $E^N$ 与 $(\mathbb{Z}/N\mathbb{Z})^2$ 之间的同构（这里 $E^N$ 是 $E$ 上 $N$ 阶点构成的群）。

考虑 $p$ 处的 Hecke 对应。它是一条曲线，记为 $H_p$，参数化五元组 $(E,\alpha,E',\alpha',\beta)$，其中 $(E,\alpha)$ 与 $(E',\alpha')$ 是 $X(N)$ 的两个点，$\beta$ 是 $E\to E'$ 的同源，且 $\ker(\beta)$ 的阶为 $p$。由于 $N$ 与 $p$ 互素，$\beta$ 诱导同构 $E^N\to E'^N$；我们要求它与两个群到 $(\mathbb{Z}/N\mathbb{Z})^2$ 的等同相容。

和前面一样，有自然映射
\[
X(N)\xleftarrow{h^{\leftarrow}}H_p\xrightarrow{h^{\rightarrow}}X(N),
\]
它们定义算子 $f\mapsto(h^{\rightarrow})_!((h^{\leftarrow})^*(f))$，与作用在模形式上的 Hecke 算子一致。

\ob{2.8}{ob:2.8}
回到几何情形，用几何语言写出尖点条件。设 $n=n_1+n_2$ 是分拆。考虑短正合列
\[
0\to M_1\to M\to M_2\to 0
\]
的同构类集合 $\Fl^n_{n_1,n_2}$，其中 $M_1$ 与 $M_2$ 分别是秩 $n_1$ 与 $n_2$ 的向量丛。

有自然映射 $p$ 与 $q$
\[
\Bun_n\xleftarrow{p}\Fl^n_{n_1,n_2}\xrightarrow{q}\Bun_{n_1}\times\Bun_{n_2},
\]
分别把上述短正合列送到 $M$ 与 $(M_1,M_2)$。

定义常项算子 $r^n_{n_1,n_2}:\Funct(\Bun_n)\to\Funct(\Bun_{n_1}\times\Bun_{n_2})$ 为
\[
f\mapsto q_!(p^*(f)).
\]

现在容易看出：函数 $f$ 是尖点的，当且仅当对所有 $n_1+n_2=n$、$n_1,n_2>0$ 都有 \mbox{$r^n_{n_1,n_2}(f)=0$}。
